% ============================================================================
% Monografia de TCC I — Sistema IoT para Monitoramento da Presença de
% Estudantes em Transporte Escolar Rural por Reconhecimento de Voz e Vídeo
%
% Aluna: Fernanda Andrade Bicalho da Cruz
% Orientador: Janier Arias García — Coorientador: Frederico Gualberto Ferreira Coelho
% Universidade Federal de Minas Gerais — Escola de Engenharia
%
% Este arquivo é mantido pela skill "monografia-tex" (.claude/skills/monografia-tex).
% Ele é a fonte única do texto da monografia — não crie cópias paralelas.
% Para compilar: envie este arquivo, "referencias.bib" e a pasta "assets_tex/"
% (com as imagens da capa) juntos para um projeto no Overleaf.
% ============================================================================

\documentclass[
    12pt,
    openright,
    twoside,
    a4paper,
    chapter=TITLE,
    section=TITLE,
    english,
    brazil
    ]{abntex2}

% ----------------------------------------------------------------------------
% Pacotes
% ----------------------------------------------------------------------------
\usepackage{lmodern}
\usepackage[T1]{fontenc}
\usepackage[utf8]{inputenc}
\usepackage{lastpage}
\usepackage{indentfirst}
\usepackage{color}
\usepackage{graphicx}
\usepackage{microtype}
\usepackage{float}

\usepackage[brazilian,hyperpageref]{backref}
\usepackage[alf]{abntex2cite}

\renewcommand{\backrefpagesname}{Citado na(s) página(s):~}
\renewcommand{\backref}{}
\renewcommand*{\backrefalt}[4]{%
	\ifcase #1
		Nenhuma citação no texto.
	\or
		Citado na página #2.
	\else
		Citado #1 vezes nas páginas #2.
	\fi}%

% ----------------------------------------------------------------------------
% Dados para capa e folha de rosto
% (a capa e a folha de rosto abaixo são montadas manualmente, não via
% \imprimircapa, para controlar a posição da logo da UFMG — mas estes campos
% continuam usados pelo abntex2 em metadados do PDF e no hypersetup)
% ----------------------------------------------------------------------------
\titulo{Sistema IoT para Monitoramento da Presença de Estudantes em Transporte Escolar Rural por Reconhecimento de Voz e Vídeo}
\autor{Fernanda Andrade Bicalho da Cruz}
\local{Belo Horizonte}
\data{2026}
\orientador{Janier Arias García}
\instituicao{Universidade Federal de Minas Gerais -- Escola de Engenharia -- Curso de Engenharia de Sistemas}
\tipotrabalho{Trabalho de Conclusão de Curso I}
\preambulo{Trabalho de Conclusão de Curso I apresentado ao Curso de Engenharia de Sistemas da Escola de Engenharia da Universidade Federal de Minas Gerais como requisito parcial para obtenção do título de Bacharel em Engenharia de Sistemas.}

\makeatletter
\hypersetup{
    pdftitle={\@title},
    pdfauthor={\@author},
    pdfsubject={\imprimirpreambulo},
    pdfcreator={LaTeX com abnTeX2},
    pdfkeywords={IoT}{reconhecimento facial}{transporte escolar}{TCC},
    colorlinks=true,
    linkcolor=blue,
    citecolor=blue,
    filecolor=magenta,
    urlcolor=blue,
    bookmarksdepth=4
}
\makeatother

\makeindex

% ============================================================================
\begin{document}

\selectlanguage{brazil}
\frenchspacing

% ----------------------------------------------------------------------------
% ELEMENTOS PRÉ-TEXTUAIS
% ----------------------------------------------------------------------------
\pretextual

% --- Capa (manual, com logo da UFMG) ---
\begin{titlepage}
\begin{center}
    \includegraphics[width=4cm]{assets_tex/ufmg_logo.png}
    \vspace{0.5cm}

    {\large UNIVERSIDADE FEDERAL DE MINAS GERAIS}\\
    {\large ESCOLA DE ENGENHARIA}\\
    {\large CURSO DE ENGENHARIA DE SISTEMAS}

    \vspace{3cm}
    {\large\textbf{\imprimirautor}}
    \vspace{3cm}

    {\LARGE\textbf{\imprimirtitulo}}

    \vfill
    {\large \imprimirlocal}\\
    {\large \imprimirdata}
\end{center}
\end{titlepage}

% --- Folha de rosto ---
\begin{titlepage}
\begin{center}
    {\large\textbf{\imprimirautor}}
\end{center}

\vspace{3cm}
\begin{center}
    {\LARGE\textbf{\imprimirtitulo}}
\end{center}
\vspace{2cm}

\begin{flushright}
    \begin{minipage}{8cm}
        \imprimirpreambulo

        \vspace{0.5cm}
        Orientador: \imprimirorientador \\
        Coorientador: Frederico Gualberto Ferreira Coelho
    \end{minipage}
\end{flushright}

\vfill
\begin{center}
    {\large \imprimirlocal}\\
    {\large \imprimirdata}
\end{center}
\end{titlepage}

% --- Resumo ---
\begin{resumo}
    % TODO(monografia-tex): substituir por um resumo de 1 parágrafo assim que
    % Introdução, Metodologia e Revisão Bibliográfica estiverem redigidas —
    % o resumo é escrito por último, resumindo o que já foi escrito, nunca o
    % contrário.

    \vspace{\onelineskip}
    \noindent
    \textbf{Palavras-chave}: IoT. reconhecimento facial. reconhecimento de voz. transporte escolar. TRANSCOLAR RURAL.
\end{resumo}

\begin{resumo}[Abstract]
    \begin{otherlanguage*}{english}
        % TODO(monografia-tex): tradução do resumo em inglês.

        \vspace{\onelineskip}
        \noindent
        \textbf{Keywords}: IoT. facial recognition. voice recognition. school transportation.
    \end{otherlanguage*}
\end{resumo}

% --- Sumário ---
\pdfbookmark[0]{\contentsname}{toc}
\tableofcontents*
\cleardoublepage

% ----------------------------------------------------------------------------
% ELEMENTOS TEXTUAIS
% ----------------------------------------------------------------------------
\textual

\chapter{Introdução}

O transporte escolar rural é essencial para garantir o acesso à educação em
regiões afastadas, mas apresenta desafios relacionados à segurança dos
estudantes, ao controle de presença e à eficiência operacional da gestão
pública. Na ausência de mecanismos automatizados, o registro de embarque e
desembarque depende de processos manuais, sujeitos a falhas e de difícil
auditoria, o que compromete tanto a segurança física dos alunos quanto a
transparência da gestão do serviço.

Este trabalho propõe a concepção e a implementação de um dispositivo de
Internet das Coisas (IoT) embarcado em veículos escolares rurais, capaz de
identificar e registrar a presença dos estudantes por meio de reconhecimento
de voz e de vídeo. Os eventos de presença capturados pelo dispositivo são
integrados de forma automatizada ao sistema TRANSCOLAR RURAL, permitindo à
gestão pública acompanhar a operação do transporte escolar com maior
precisão e agilidade.

O objetivo geral deste trabalho é projetar e implementar um dispositivo IoT
embarcado em ônibus escolares que registre a presença dos alunos por
reconhecimento de voz e vídeo, integrando os dados gerados ao sistema
TRANSCOLAR RURAL. A partir desse objetivo geral, definem-se os seguintes
objetivos específicos: (i) desenvolver um protótipo de hardware IoT dotado de
câmera e microfone; (ii) implementar algoritmos de reconhecimento facial e de
voz adequados às condições reais de operação; (iii) integrar o dispositivo ao
sistema TRANSCOLAR RURAL por meio de rede móvel; (iv) viabilizar a geração de
relatórios inteligentes voltados a gestores escolares; e (v) avaliar o
desempenho, a segurança e a confiabilidade do sistema proposto.

O desenvolvimento deste trabalho está organizado em dois ciclos. O TCC~I, ao
qual se refere esta monografia, compreende o levantamento de requisitos, a
fundamentação teórica sobre IoT e Inteligência Artificial aplicadas ao
problema, a definição da arquitetura de hardware e software do sistema e a
avaliação dos impactos socioeconômicos, legais e ambientais da tecnologia
proposta. O TCC~II, por sua vez, dará continuidade ao trabalho com o
desenvolvimento prático do software embarcado, a integração de comunicação
via rede móvel, a simulação em ambiente controlado, a realização de testes
piloto em ônibus escolares rurais e a validação sistêmica de métricas de
acurácia e tempo de resposta.

Esta monografia está organizada da seguinte forma: o Capítulo~2 apresenta a
revisão bibliográfica sobre reconhecimento biométrico e comunicação para IoT;
o Capítulo~3 descreve a metodologia adotada para alcançar os objetivos do
TCC~I; o Capítulo~4 discute os impactos socioeconômicos, legais e de
humanidades da tecnologia proposta; e o Capítulo~5 apresenta o cronograma
planejado para o TCC~II.

\chapter{Revisão Bibliográfica}
% TODO(monografia-tex): ainda não escrita. Deve cobrir: técnicas de
% reconhecimento facial (OpenCV/algoritmos candidatos), reconhecimento de voz
% em ambientes ruidosos, protocolos de comunicação leve para IoT (MQTT).

\chapter{Metodologia}
% TODO(monografia-tex): ainda não escrita. Deve descrever como os objetivos de
% TCC I serão alcançados (levantamento de requisitos, revisão bibliográfica,
% projeto arquitetural e análise de trade-offs, avaliação de impactos).

\chapter{Avaliação de Impactos Socioeconômicos, Legais e de Humanidades}
% TODO(monografia-tex): ainda não escrita. Deve abordar segurança física dos
% estudantes, transparência/otimização de recursos públicos, restrições legais
% de privacidade e proteção de dados de menores (LGPD), e impacto ambiental do
% descarte de hardware.

\chapter{Cronograma do TCC II}
% TODO(monografia-tex): ainda não escrita. Deve detalhar as atividades
% previstas para o TCC II (desenvolvimento do software embarcado, integração
% via MQTT/API, simulação, teste piloto, validação de métricas).

% ----------------------------------------------------------------------------
% ELEMENTOS PÓS-TEXTUAIS
% ----------------------------------------------------------------------------
\postextual

\bibliography{referencias}

\end{document}
